% Generated draft: requires mathematical and source review.
\documentclass[11pt]{article}
\usepackage[margin=25mm]{geometry}
\usepackage{fontspec}
\setmainfont{Latin Modern Roman}
\usepackage{amsmath,amssymb,mathtools,unicode-math}
\setmathfont{Latin Modern Math}
\usepackage{enumitem,needspace,xurl,hyperref}
\setlength{\parindent}{0pt}
\setlength{\parskip}{5pt}
\setlength{\emergencystretch}{3em}
\newcommand{\N}{\mathbb{N}}
\newcommand{\Z}{\mathbb{Z}}
\newcommand{\Q}{\mathbb{Q}}
\newcommand{\R}{\mathbb{R}}
\newcommand{\C}{\mathbb{C}}
\newcommand{\F}{\mathbb{F}}
\newcommand{\A}{\mathbb{A}}
\newcommand{\PP}{\mathbb{P}}
\newcommand{\E}{\mathbb{E}}
\newcommand{\eps}{\varepsilon}
\newcommand{\dd}{\,\mathrm{d}}
\begin{document}
\section*{UnsolvedMath ID: 6800016}
\textbf{Source problem code:} catalogue entry
\subsection{Definitions and mathematical statement}
In the usual Euclidean interpretation, let \(\Sigma\subset\mathbb R^3\) be a smooth complete embedded minimal surface, without boundary, of finite topological genus. For a two-sided surface, its normal second-variation form is
\[
 Q_\Sigma(u,u)=\int_\Sigma\bigl(|\nabla u|^2-|A|^2u^2\bigr)\,\dd A,
 \qquad u\in C_c^\infty(\Sigma),
\]
where \(A\) is the second fundamental form. Its Morse index is the supremum of dimensions of subspaces on which \(Q_\Sigma\) is negative definite. The source asks whether there exists an embedded finite-genus minimal surface with
\[
 \operatorname{ind}(\Sigma)=4.
\]
The excerpt does not actually specify ambient space, completeness, boundary or one-/two-sided conventions; the Euclidean complete-surface formulation is explicitly a standard contextual reading to be checked, not a restriction silently supplied as original wording. No finite-total-curvature or prescribed-end hypothesis is added.

The more focused question concerns the source's one-parameter deformation \((\Sigma_t)_{t\in I}\) of Costa's surface: is
\[
 \operatorname{ind}(\Sigma_s)=\operatorname{ind}(\Sigma_t)
 \qquad\text{for every }s,t\in I?
\]
No parameter interval or construction of this family is supplied, so it cannot be identified with a guessed deformation family or endpoint degeneration. The constancy question does not prescribe the common numerical index.

The research summary repeatedly asserts an index of \(3\) for Costa's surface and discusses possible jumps to \(4\), but supplies no supporting index theorem. That numerical assertion is retained only as an unverified source claim and is not used to change the actual questions. The summary also contains the unexplained word ``clavicles'' and mentions Maximo--Nunes work without a precise reference. Its partial/uncertain assessment reports no verified definitive resolution; no current status or numerical index is independently certified here.
\subsection{Short English statement}
Does an embedded minimal surface of finite genus have Morse index four, and is the index constant along the intended one-parameter Costa deformation? The ambient/completeness setting and the family are missing, and the summary's unsupported numerical Costa-index claim is not imposed on the questions.
\subsection{Sources}
\begin{itemize}
\item Local registry: \url{data/unsolvedmath/problems.json}, permanent UnsolvedMath ID \(6800016\), source code catalogue entry; supplied statement, background, and research summary, updated 2026-08-06. Proposed by Davi Maximo, University of Pennsylvania, 2018.
\item Morgan and Pansu, \emph{A List of Open Problems in Differential Geometry} (2018), Question environment 16, ``Morse index of embedded minimal surfaces'': \url{https://www.imo.universite-paris-saclay.fr/~pansu/problems_MTDG.tex}. Supplied source-TeX extraction accessed 2026-07-29.
\item Costa's surface/deformation and Maximo--Nunes work are mentioned without a family definition, exact index citation, or full publication details.
\end{itemize}
\end{document}
